\section{FunMentor}
\begin{wrapfigure}{r}{0.5\textwidth}
    \centering
    \vspace{-30pt}
    \includegraphics[width=\linewidth]{figures/Funmentor.pdf}
    \vspace{-20pt}
    \caption{\textbf{FunMentor's Refining Process.} FunMentor asks multi-round questions to help VLMs to generate persuading answers.}
    \vspace{-30pt}
    \label{fig:funmentor}
    \vspace{10pt}
\end{wrapfigure}
This section presents the details of \textbf{FunMentor} for counter-intuitiveness understanding. FunMentor is an agent that refines a VLM's answer through multi-turn dialogues, ensuring it generates answers that best explain the given surprising content. The refining process comprises three components:\\
\noindent\textbf{Real Fact Collection}\quad 
Initially, FunMentor poses a series of inquiries to the VLM model it aims to assist, focusing on fact-aware questions to accurately comprehend the objective content of the video (e.g., FunMentor might ask: ``Please describe the items and characters appearing in each frame of the video''). This step is crucial because, without this preliminary context, FunMentor would have no knowledge of the video content and might be gullible to VLM's humorous explanations of the video. Thus, providing it with some factual clues is very important.\\
\noindent\textbf{Answer Judgement}\quad 
FunMentor assesses VLMs' responses based on several key aspects: \textbf{1) Language Fluency}, ensuring responses are grammatically fluent; \textbf{2) Detail Accuracy}, verifying factual correctness in relation to the video's content; \textbf{3) Logical Reasoning}, checking for coherent logic and smooth transitions; and \textbf{4) Humorous Insight}, expecting responses to provide humor beyond a mere description of the video. If not passed, FunMentor will try to generate some feedback, which is explained below.\\
\noindent\textbf{Suggestion Generation}\quad 
Upon receiving unsatisfactory responses, FunMentor formulates constructive suggestions and tailored prompts, which are then sent back to the VLMs for revision, anticipating their improved answers. This process involves the agent analyzing the model's initial response, blending pre-defined prompts with the context of the original question, to generate specific feedback instructions. For example, FunMentor examines the issues in the VLM's response and advises on what should and shouldn't be included in the revised answer, guiding the VLM towards a more accurate and relevant reply.\\
Fig. \ref{fig:funmentor} shows the pipeline of FunMentor. More details as shown in Appendix \ref{appendix_c: More Details of FunMentor}.

% \noindent\textbf{Quality Inspection of Responses.}\quad 
% The agent first evaluates whether the model's answers genuinely comprehend humor, specifically, xxxx. Throughout the dialogue, the agent continuously assesses the quality of responses until the model produces a satisfactory answer.\\
% \noindent\textbf{Prompt Generation}\quad 
% When the model provides subpar responses, the agent combines pre-defined prompts, the model's response, and the context of the VQA to generate feedback instructions for the model.\\
% \noindent\textbf{VLM Fine-tuning}\quad 
% The model undergoes end-to-end training, using the provided prompts to generate revised answers.\\
% Throughout the entire process, FunMentor functions as a humor gatekeeper, controlling the quality of responses provided by VLM through promtpts until the answer is deemed satisfactory for understanding humor. FunMentor adapts well to various VLMs and can extend its capabilities beyond HumorQA to comprehend other counter-intuitive video contexts. More details can be seen in Appendix \ref{appendix_d: Exp_result}. The following experiments showcase FunMentor's results with VLMs.
\section{Experiments}
% \jingkang{Please have 1-2 sentences to introduce what will you show in this section.}
In this section, we begin by introducing various caption-based and instruction-based models for evaluation. We then delve into diverse metrics tailored for FunQA tasks, with extensive details provided in Appendix \ref{appendix_d: Experiments}. Our comprehensive experiments and deep analysis of the results are then presented. Additionally, we conduct a comparative study contrasting FunQA with the previous VQA dataset, NExT-QA, which similarly features multiple-choice and open-ended questions. This comparison underscores the unique attributes and importance of FunQA.

% \subsection{Baselines}
% \textbf{1) Q-ONLY} This baseline only uses the question feature as input to predict the answer.
% \textbf{2) V-ONLY} This baseline only uses the visual feature as input to predict the answer.
% \subsubsection{Caption-based Models}

% \noindent\textbf{mPLUG}\quad
% mPLUG \cite{li2022mplug} consists of two unimodal encoders for image and text independently, a cross-modal skip-connected network, and a decoder for text generation. Based on the connected representation of the image and prefix sub-sequence, the decoder is trained with a prefix language modeling loss by generating the remaining caption. \\
% \noindent\textbf{GIT}\quad
% GIT \cite{wang2022git} is composed of one image encoder and one text decoder, for videos, multiple frames are sampled and encoded independently, and features are added with an extra learnable temporal embedding before concatenation. The image encoder is based on the contrastive pre-trained model \cite{yuan2021florence} with raw image input and a compact 2D feature map output. The text decoder is a transformer module to predict the text description.
%which consists of multiple transformer blocks, each of which is composed of one self-attention layer and one feed-forward layer.


% \subsubsection{Instruction-based Models}
% \noindent\textbf{VideoChat}\quad
% VideoChat \cite{2023videochat} is an end-to-end chat-centric video understanding system, the version we use is VideoChat-13B. Its VideoChat-Embed architecture is instantiated using BLIP-2 \cite{li2023blip2} and StableVicuna (13b-delta), and combines pre-trained ViT-G \cite{sun2023evaclip} and GMHRA \cite{wang2022internvideo}. For the token interface, VideoChat-Embed uses the pre-trained QFormer with an additional linear projection to output Video Embedding. Both video description and Video Embedding will be input in LLMs (LLAMA-13B \cite{touvron2023llama}) for multimodal understanding and output timestamped video text descriptions. 
%which has a VideoChat-Text architecture with video keyframes extracted using ffmpeg as input and a timestamped video text description as output
% \\
% \noindent\textbf{Video-ChatGPT}\quad 
% Video-ChatGPT \cite{maaz2023videochatgpt} is a large vision-language model with a dedicated video-encoder and LLM, which feeds the video frames into pre-trained video encoder, adds spatio-temporal features and feeds them into linear layer. Video Embedding is then input into LLM (Vicuna-7B, v1.1) along with System command and User Query to output the answer. It uses a data-centric, human-assisted, and semi-automated annotation framework for high-quality video instructional data with unique multimodal (visual-verbal) capabilities.\\
% \noindent\textbf{Otter}\quad 
% The Otter model \cite{li2023otter} utilizes the OpenFlamingo training paradigm, where the pretrained OpenFlamingo model consists of a LLaMA-7B language encoder \cite{touvron2023llama} and a CLIP ViT-L/14 vision encoder \cite{radford2021learning}. In the fine-tuning process for instruction tuning, Otter freezes both the encoders and only fine-tunes the Perceiver resampler module. We evaluate two versions of Otter: one that is fine-tuned on the Dense Caption dataset \cite{krishna2017dense}, and another that is fine-tuned on the FunQA training set. Due to GPU memory limitations, Otter can only train and test on 128 frames from a video.

\subsection{Model Zoo}
We categorize the models capable of addressing the majority of VQA tasks into two classes: Caption-based models and Instruction-based models. Caption-based models generate suitable captions for videos by taking different prompts as input. For this category, we evaluate mPLUG~\cite{li2022mplug} and GIT \cite{wang2022git}. Meanwhile, Instruction-based models are capable of answering a wide array of questions. For this category, we evaluate models VideoChat \cite{2023videochat}, Video-ChatGPT \cite{maaz2023videochatgpt}, mPLUG-Owl \cite{ye2023mplugowl}, Video-LLaMA \cite{damonlpsg2023videollama} and Otter \cite{li2023otter}.
For Otter, we evaluate two versions: one is fine-tuned on the Dense Caption \cite{krishna2017dense}, and another is fine-tuned on the FunQA training set. We evaluate the proposed FunMentor based on Video-ChatGPT and Otter.

% \textbf{Moment} 
% We use intersection of union to evaluate moment.\\
% \textbf{Score} The metric is $ (result - score)/20$\\
% \textbf{Free-text}
% 1) NLG. NLG traditional metrics such as BLUE \cite{papineni2002bleu}, ROUGE-1/2/L \cite{lin2004rouge} and METEOR \cite{banerjee2005meteor} are common for natural language understanding, which enables researchers to fairly compare their models to others. In GEM \cite{gehrmann2021gem}, people . However, most of the tasks in our dataset are free-text, which is different from the rest of the extant VQA datasets, and most of the existing models fail to yield the desired results, which are close to 0 when evaluated with the existing NLG metric after FunQA testing.\\
% 2) GPT 
% Some work \cite{} has demonstrated that GPT can be used as a metric for NLG with good results, so we introduced GPT to help with the free-text evaluation. We carefully designed the prompt to achieve the desired effect, and GPT will focus on the consistency of the details of the two texts and the understanding of humor, creativity and magic effects, as detailed in \ref{Metics}.
\subsection{Evaluation Metrics}
\noindent\textbf{Timestamp Localization (H1, C1, M1)}\quad 
We employ the intersection of unions (IOU) based on time span. \\
% which is calculated as $\min(result\_end, groud\_truth\_end) - \max(result\_start, groud\_truth\_start)/\min(result\_end, groud\_truth\_end) - \max(result\_start, groud\_truth\_start)$ 
\noindent\textbf{Description \& Reasoning (H2-4, C2-4, M2-3)}\quad
For all the \textbf{free-text} tasks, we employ two approaches for evaluation. Firstly, we utilize traditional NLG (Natural Language Generation) metrics. We use BLEU-4 \cite{papineni2002bleu}, ROUGE-L \cite{lin2004rouge}, CIDEr \cite{vedantam2015cider}, and BLEURT \cite{pu2021learning} as our metrics. The first two rely on N-gram overlap, which is only sensitive to lexical variations and cannot identify changes in sentence semantics or grammar. The latter two are reference-based evaluation metrics. Secondly, several works \cite{wang2023chatgpt, fu2023gptscore, kamalloo2023evaluating, chiang2023can} have shown promising results in utilizing GPT as a metric for NLG. Therefore, we introduce GPT-4 to assist in evaluating free-text similarity. We carefully design the prompts to make it possible to give objective ratings as much as possible like a human being. More details of GPT-4 prompts and evaluation criteria are provided in Appendix 
\ref{appendix_d: Evaluation Metrics}.\\
\noindent\textbf{Creative Scoring (C5)}\quad $Creative Score Metric = 100 \times \left(1 - \frac{{\lvert Predict - GT \rvert}}{{20}}\right)$.
\subsection{Results and Observations}
\begin{table}[t]
\caption{\textbf{Main Results on FunQA Benchmark.} The FunQA benchmark consists of four task categories. H1, C1, M1 represent the counter-intuitive timestamp localization task, where \textbf{IOU} is used as the metric. H2, C2, M2 represent the detailed video description task, and H3, C3, M3 represent reasoning around counter-intuitiveness. For the higher-level tasks, H4, C4 involve attributing a fitting and vivid title. The responses for all these tasks in free-text format. We utilize the \textbf{BLEURT} (B.) and \textbf{GPT-4} metrics for evaluation. The scores for the traditional metrics (BLEU-4 and CIDEr) are all close to zero. Here, we present only the BLEURT and GPT-4 scores, with the full results available in the Appendix \ref{appendix_d: Exp_result}). C5 represents scoring the video creativity, and the metric is the \textbf{Accuracy} between the predicted score and the official score. We tested the caption-based and instruction-based models. 
Here we clarify the abbreviation in the table.  \textbf{L.M.}: GIT\_LARGE\_MSRVTT; \textbf{L.V.}: GIT\_LARGE\_VATEX; \textbf{D.C.} means finetuned on Dense Caption; \textbf{FunQA} means finetuned on FunQA. }
\label{T:experiment}
\vspace{-10pt}
{\renewcommand\baselinestretch{1.5}\selectfont
% \Large
\setlength\tabcolsep{10pt} 
\resizebox{1.0\linewidth}{!}{
\begin{tabular}{lp{1cm}<{\centering} cccp{1cm}<{\centering}ccccp{1cm}<{\centering}cc}
\toprule
\multirow{2}{*}{} & \multicolumn{4}{c}{HumorQA} & \multicolumn{5}{c}{CreativeQA} & \multicolumn{3}{c}{MagicQA} \\

\cmidrule(l){2-5} \cmidrule(l){6-10} \cmidrule(l){11-13}
Task & H1 & H2-Des. & H3-Rea. & H4-Title & C1 & C2-Des. & C3-Rea. & C4-Title & C5-Score & M1 & M2-Des. & M3-Rea. \\
\toprule
Metrics & IOU & B. / GPT-4 & B. / GPT-4 & B. / GPT-4 & IOU & B. / GPT-4 & B. / GPT-4 & B. / GPT-4 & Acc & IOU & B. / GPT-4 & B. / GPT-4 \\
% \toprule
% \rowcolor{TRAD_EVAL}
\toprule
\rowcolor{C-Model}
\multicolumn{13}{l}{\textbf{- Caption-based Model}} \\
\rowcolor{C-Model}
mPLUG \cite{li2022mplug}  
& 0.0
& 19.9 / 3.9
& \textbf{25.7} / \textbf{6.0}
& 22.1 / \textbf{11.2}

& 0.0 
& 14.9 / 3.0
& \textbf{24.2} / \textbf{6.9}
& \textbf{20.8} / \textbf{18.8}
& 0.0 / 0.0 

& 0.0 
& 19.7 / 4.0
& \textbf{21.2} / \textbf{8.1} \\
\rowcolor{C-Model}
GIT (L.M.) \cite{wang2022git} 
& 0.0 
& 22.4 / 3.6
& 0.0 / 0.0
& 17.0 / 8.9

& 0.0 
& 14.4 / 3.8
& 0.0 / 0.0
& 7.1 / 11.3
& 0.0 / 0.0

& 0.0
& 19.4 / 8.2
& 0.0 / 0.0 \\

\rowcolor{C-Model}
GIT (L.V.) \cite{wang2022git} 
& 0.0 
& \textbf{33.3} / \textbf{4.0}
& 0.0 / 0.0
& \textbf{25.9} / 10.0

& 0.0 
& \textbf{20.5} / \textbf{4.2}
& 0.0 / 0.0
& 10.5 / 12.0
& 0.0 / 0.0 

& 0.0 
& \textbf{29.8} / \textbf{8.6}
& 0.0 / 0.0 \\
\midrule
\rowcolor{I-Model}
\multicolumn{13}{l}{\textbf{- Instruction-based Model}} \\
\rowcolor{I-Model}
VideoChat \cite{2023videochat} 
& 0.0 
& 44.0 / 17.9
& 45.4 / 31.9
& 20.2 / 31.7

& 0.0 
& 21.7 / 5.9
& 22.8 / \textbf{17.7}
& 7.3 / 31.1
& 67.5 

& 0.0 
& 47.4 / 8.2
& 43.1 / 44.6 \\


\rowcolor{I-Model}
Video-ChatGPT \cite{maaz2023videochatgpt} 
& 0.0 
& 39.9 / \textbf{24.3}
& 40.1 / 24.9
& 36.5 / 41.2

& 0.0 
& \textbf{45.8} / 6.6
& 45.2 / 9.1
& 30.9 / \textbf{48.8}
& \textbf{85.4} 

& 0.0 
& \textbf{50.8} / \textbf{11.2}
& 43.3 / 40.4 \\


\rowcolor{I-Model}
mPLUG-Owl \cite{ye2023mplugowl} 
& 0.0 
& 44.5 / 10,7
& \textbf{47.3} / \textbf{35.0}
& 29.8 / \textbf{48.8}

& 0.0 
& 43.0 / 5.0
& \textbf{44.7 }/ 10.6
& 23.9 / 36.3
& 66.7 

& 0.0 
& 46.4 / 8.6
& 43.9 / 30.9 \\

\rowcolor{I-Model}
Video-LLaMA \cite{damonlpsg2023videollama} 
& 0.0 
& \textbf{48.4} / 7.7
& 42.9 / 29.0
& 46.5 / 34.1

& 0.0 
& 45.5 / 7.2
& 41.1 / 17.2
& 42.3 / 31.2
& 64.2 

& 0.0 
& 50.1 / 10.2
& 39.0 / 28.0 \\

\rowcolor{I-Model}
Otter (D.C.) \cite{li2023otter} 
& 0.0 
& 30.2 / 7.7
& 32.3 / 28.3
& 21.7 / 20.0

& 0.0 
& 28.7 / 1.7
& 32.9 / 7.9
& 17.7 / 36.3
& 45.0

& 0.0 
& 32.5 / 2.1
& 27.3 / 36.8 \\
\rowcolor{I-Model}
Otter (FunQA) \cite{li2023otter} 
& 0.0 
& 38.4 / 8.9
& 42.6 / 31.7
& \textbf{47.5} / 32.1

& 0.0 
& 40.0 / \textbf{7.3}
& 41.1 / 8.8
& \textbf{44.5} / 38.8
& 69.4

& 0.0 
& 44.7 / 10.3
& \textbf{44.5} / \textbf{47.5} \\
\midrule
\rowcolor{O-Model}
Video-ChatGPT \cite{maaz2023videochatgpt}  + \textbf{FunMentor (Ours)}
& 0.0 
& \textbf{65.2} / \textbf{33.2}
& \textbf{57.5} / 36.5
& 50.2 / \textbf{65.1}

& 0.0 
& \textbf{66.3} / \textbf{14.2}
& \textbf{58.7} / \textbf{23.4}
& 45.3 / \textbf{52.2}
& \textbf{85.4}

& 0.0
& \textbf{55.1} / \textbf{13.3}
& \textbf{46.3} / 54.8 \\
\rowcolor{O-Model}
Otter (FunQA) \cite{li2023otter} + \textbf{FunMentor (Ours)}
& 0.0
& 33.4 / 13.4
& 37.8 / \textbf{45.8}
& \textbf{58.3} / 34.2

& 0.0 
& 60.4 / 11.0
& 44.4 / 9.3
& \textbf{53.9} / 43.5
& 69.4

& 0.0 
& 43.5 / 12.81
& 38.91 / \textbf{56.4}\\



\bottomrule

\end{tabular}}\par
}
\vspace{-0.5cm}
\end{table}

In Table \ref{T:experiment}, we show the results of model zoo and
our proposed method. For clarity, we provided a list of dos and don’ts for comparing values in the table:\\ 
\textbf{a. Values from different tasks with the same video type (e.g., H2 and H3) are not comparable.} We observe that the model output in reasoning tasks may contain several words that match the ground truth (GT) while the meaning might be incorrect, resulting in inflated results. Therefore, our comparison is only based on qualitative analysis. \\
\textbf{b. Same metric with different models (see vertically) is comparable.}\\
\textbf{c. Same task for different video types (e.g., H2 and C2) is comparable.}\\
As an example, Fig. \ref{fig:vlm's response} illustrates the responses of VLM before and after applying our method to the HumorQA dataset. Overall, the performance of the models on the FunQA dataset is generally unsatisfactory. However, after fine-tuning VLM with FunMentor, a notable improvement is shown in its ability to comprehend counter-intuitiveness. We have made several key findings:\\
\noindent\textbf{Timestamp localization task is the most challenging.}\quad
Caption-based models focus mainly on captioning and often omit temporal information. 
\begin{wraptable}{r}{5.0cm}
\centering
\vspace{-1cm}
\small
\caption{\textbf{Timestamp Localization Task's baseline.} The result uses IOU metrics.}
\resizebox{\linewidth}{!}{
\begin{tabular}{lp{2cm}<{\centering}p{2cm}<{\centering}p{2cm}<{\centering}}
\toprule
Task         & H1                              & C1                             & M1                             \\ \midrule
Random   &  29.4 & 32.6 &  21.1   \\
Mean  &  55.5  &  82.1 &  46.3 \\
\midrule
TimeChat \cite{Ren2023TimeChat} &  8.4  &  14.6 &  6.0   \\
\bottomrule 
\end{tabular}}
\label{tab:iou}
\caption{\textbf{Human Performance on HumorQA.} The average scores using BLEURT / GPT-4. Full results are in Appendix \ref{appendix_d: Experiments}.}
\resizebox{\linewidth}{!}{
\begin{tabular}{lccc}
\toprule
Task         & H2                              & H3                             & H4                              \\ \midrule
Model SOTA   &   65.2 / 33.2  &   57.5 / 45.8 &  58.3 / 65.1   \\
Human &  70.0 / \textbf{83.6} &  63.1 / \textbf{84.4} &  67.3 / \textbf{79.1}\\
\bottomrule 
\end{tabular}}
\vspace{-1cm}
\label{tab:human}
\end{wraptable}
Due to this, they are not equipped to deal with our task that emphasizes temporal aspects, so they are not scored. Similarly, Instruction-based models, such as Otter, take visual information from specific frames without temporal context. Their outputs are thus confined to individual frames, making them ineffective at addressing temporal localization. In summary, none of the current VLMs could solve H1, C1, or M1 tasks since they do not have a sense of time (refer to Appendix \ref{appendix_d: Failure of t1}). To enhance the timestamp localization task analysis, we add comparative baselines with Random Guess; Mean Time Span derived from the training dataset for an average reference; and TimeChat \cite{Ren2023TimeChat}, a specialized time-sensitive multimodal language model. As shown in Table \ref{tab:iou}, TimeChat’s score is the lowest showing the challenges of the Timestamp Localization Task.\\
% \begin{center}
% \begin{minipage}{\textwidth}
% \begin{minipage}[t]{0.45\textwidth}
% \makeatletter\def\@captype{table}
% \caption{\textbf{Timestamp Localization Task's baseline.} This table shows the Timestamp Localization Task's baseline with Random Guess; Mean Time Span derived from the training dataset for an average reference; and Timechat \cite{Ren2023TimeChat}, a specialized time-sensitive multimodal language model.}
% \vspace{-10pt}
% \resizebox{0.8\linewidth}{!}{
% \begin{tabular}{lccc}
% \toprule
% Task         & H1                              & C1                             & M1                             \\ \midrule
% Random   &  29.4 & 32.6 &  21.1   \\
% Mean  &  55.5  &  82.1 &  46.3 \\
% \midrule
% TimeChat \cite{Ren2023TimeChat} &  8.4  &  14.6 &  6.0   \\
% \bottomrule 
% \end{tabular}}
% \label{tab:iou}
% \end{minipage}
% \begin{minipage}[t]{0.55\textwidth}
% \makeatletter\def\@captype{table}
% \centering
% \caption{\textbf{Human Evaluations on HumorQA.} The tableshows average scores using BLEURT / GPT-4. Full results are in Appendix ??.}
% \label{T:human_evaluation}
% \vspace{-10pt}
% \centering
% %\setlength\tabcolsep{15pt}
% \resizebox{\linewidth}{!}{
% \begin{tabular}{lccc}
% \toprule
% Task         & H2                              & H3                             & H4                              \\ \midrule
% Model SOTA   &   65.2 / 33.2  &   57.5 / 45.8 &  58.3 / 65.1   \\
% Human &  70.0 / \textbf{83.6} &  63.1 / \textbf{84.4} &  67.3 / \textbf{79.1}\\
% \bottomrule 
% \end{tabular}
% }
% \end{minipage}
% \end{minipage}
% \end{center}
\begin{figure}[t]
    \centering
    \vspace{-10pt}
    \includegraphics[width=0.9\linewidth]{figures/model_comparsion.pdf}
    \vspace{-10pt}
    \caption{\textbf{VLM responses on before and after FunMentor.} Here shows the answers given by Video-ChatGPT~\cite{maaz2023videochatgpt} on HumorQA video before and after FunMentor. From the GroundTruth, it is evident that \textbf{three} key elements contribute to the humor in the video: \textbf{``sock''}, \textbf{``mask''}, and \textbf{``cat resemble a human''}. Initially, Video-ChatGPT only identified the mask, failing to grasp the full essence of the video's humor. However, after the combination with FunMentor, Video-ChatGPT successfully recognized the \textbf{``looks like a person''} and the \textbf{``fact that it is wearing a sock''}, thus demonstrating a true understanding of what makes the video funny.
    % On task \textbf{H2}, \textbf{H3}, VideoChat has the best performance. On task \textbf{H4}, Video-ChatGPT and Otter answer better, which is in line with our experiment result in Table \ref{T:experiment}. However, the answers from all models are still far from the ground truth. The descriptions of details and counter-intuitive explanations have numerous shortcomings. For example, Video-ChatGPT added incorrect details to the description, such as "wearing sunglasses", the humorous reason for "throwing ketchup" was wrongly interpreted by VideoChat as "knocking over the ketchup bottle", etc. 
    }
    \label{fig:vlm's response}
    \vspace{-25pt}
\end{figure}\noindent\textbf{No clear winner across all tasks.}\quad
Caption-based models excel in providing detailed descriptions but struggle in tasks that require reasoning, resulting in a notable performance gap between description tasks (e.g., H2) and reasoning tasks (e.g., H3). On the other hand, instruction-based models demonstrate stronger reasoning capabilities than caption-based models but tend to underperform in description tasks. 
One possible explanation is that instruction-based models may generate excessive information in their answers, including a significant amount of incorrect information. We conducted experiments to compare machine versus human performance on FunQA as shown in Table \ref{tab:human}.\\
\noindent\textbf{Performance varies greatly across different video types.}\quad
% Most models can get a better answer in humor and magic videos but have difficulty answering the questions in creative videos. 
Generally CreativeQA is the most challenging especially in reasoning. For instance, the GPT-4 scores of Video-ChatGPT and Otter on C3-Reasoning are notably low. One possible reason is that humor and magic videos often depict daily life that models have encountered previously, whereas creative videos that models have never seen before, causing them unable to understand and generate reasonable answers.\\
\noindent\textbf{Insufficient evaluation metrics for free-text tasks.}\quad
Traditional metrics yield near-zero scores on free-text questions (refer to the complete FunQA Benchmark results in the Appendix \ref{appendix_d: Experiments}), as they solely focus on short textual similarity. While BLEURT scores are significantly higher, they still fall short in evaluating more complex similarities. Intuitively, GPT-4 is found to show preliminary capabilities in assessing free-text in deep understanding, which will be detailed in Appendix \ref{appendix_d: Evaluation Metrics}. However, there are still issues of instability, where the same content can receive different scores.\\
\noindent\textbf{Finetuned Otter performs well yet with limitations.}\quad
We finetuned Otter on Dense Caption and FunQA, and Otter (FunQA) shows obvious performance advantages over Otter (D.C.). As shown in Table \ref{T:experiment}, Otter (FunQA) performs better in BLEURT and GPT-4. However, there are still some limitations that Otter (FunQA) falls short of reaching SOTA scores. One possible reason revealed is that the input of Otter is only 128 frames sampled from the video, which is insufficient for comprehensive reasoning.\\
\noindent\textbf{FunMentor demonstrates the powerful potential of agent-based finetuning}\quad Our proposed FunMentor outperforms the previous best method by 10.8 for the GPT-4 score of H3 task, with improvements of 4.8 and 20.9 for Otter (FunQA) and Video-ChatGPT, respectively. Additionally, the results reveal that FunMentor achieves significant performance improvements for Video-ChatGPT, particularly in the H2 and H4 tasks, while the improvement for H3 is relatively modest. This indicates that counter-intuitive reasoning remains a challenging aspect. The substantial performance enhancement by FunMentor highlights the promising prospects of agent-based fine-tuning methods. In the context of VLM requiring extensive training data, research in this direction holds the potential to uncover the vast capabilities of VLM.

\subsection{Comparison with Previous Benchmarks}
\label{sec:compare_nextqa}
We chose NExT-QA as our comparative benchmark, which is designed to emphasize model reasoning abilities. NExT-QA also provides a multi-choice version of the dataset (NExT-OE), similar to FunQA-MC.
\begin{table}[h]
\vspace{-20pt}
\caption{\textbf{Model's Performance in NExT-OE.} This table illustrates that while the previous SOTA model, HGA, excels in the classic WUPS metric, it underperforms with GPT-4. See Appendix~\ref{appendix_d: Evaluation Metrics} for a discussion on potential issues with WUPS.}
\vspace{-10pt}
\label{T:vlm_nextoe}
\centering
\setlength\tabcolsep{5pt}
\resizebox{1.0\linewidth}{!}{
\begin{tabular}{lp{5cm}<{\centering}p{5cm}<{\centering}}
\toprule
Metrics & WUPS & GPT-4  \\
% \toprule
\hline
HGA (SOTA on NExT-OE)& 25.18	& 25.06 \\
Otter (D.C.) & 1.26 & 64.89 \\
Otter (FunQA) & 0.79 & \textbf{73.06} \\
\bottomrule
\end{tabular}
}
\vspace{-20pt}
\end{table}



\noindent\textbf{Traditional metrics are ineffective on VLMs' response.}\quad
Table \ref{T:vlm_nextoe} shows the performance on NExT-OE of Otter and HGA \cite{jiang2020reasoning} using different metrics. We evaluated these models using both traditional evaluation metrics (specifically, WUPS \cite{malinowski2014multi} that is employed by NExT-OE) and our novel evaluation metrics based on GPT-4. Otter exhibits a notably low performance on the WUPS metric that is drastically different than the GPT-4 score, primarily because WUPS is ill-suited for evaluating sentence-based responses and fares poorly when assessing phrases. Specific examples can be found in Appendix \ref{appendix_d: Exp_result}. 
% \begin{center}
% \begin{minipage}{\textwidth}

% \begin{minipage}[t]{0.48\textwidth}
% \makeatletter\def\@captype{table}
% \vspace{-0.05mm}
% \caption{\textbf{Model's Performance in NExT-OE.} This table illustrates that while the previous state-of-the-art model, HGA, excels in the classic WUPS metric, it underperforms with GPT-4. See Appendix~\ref{appendix_d: Evaluation Metrics} for a discussion on potential issues with WUPS.}
% \label{T:vlm_nextoe}
% \vspace{-5pt}
% \centering
% \setlength\tabcolsep{15pt}
% \resizebox{\linewidth}{!}{
% \begin{tabular}{l|cc}
% \toprule
% Metrics & WUPS & GPT-4  \\
% % \toprule
% \hline
% HGA (SOTA on NExT-OE)& 25.18	& 25.06 \\
% Otter (D.C.) & 1.26 & 64.89 \\
% Otter (FunQA) & 0.79 & \textbf{73.06} \\
% \bottomrule
% \end{tabular}
% }
% \end{minipage}

% \begin{minipage}[t]{0.48\textwidth}
% \makeatletter\def\@captype{table}
% \vspace{-0.05mm}
% \caption{\textbf{Model's Performance in NExT-OE.} This table illustrates that while the previous state-of-the-art model, HGA, excels in the classic WUPS metric, it underperforms with GPT-4. See Appendix~\ref{appendix_d: Evaluation Metrics} for a discussion on potential issues with WUPS.}
% \label{T:vlm_nextoe}
% \vspace{-5pt}
% \centering
% \setlength\tabcolsep{15pt}
% \resizebox{\linewidth}{!}{
% \begin{tabular}{l|cc}
% \toprule
% Metrics & WUPS & GPT-4  \\
% % \toprule
% \hline
% HGA (SOTA on NExT-OE)& 25.18	& 25.06 \\
% Otter (D.C.) & 1.26 & 64.89 \\
% Otter (FunQA) & 0.79 & \textbf{73.06} \\
% \bottomrule
% \end{tabular}
% }
% \end{minipage}
% \begin{minipage}[t]{0.48\textwidth}
% \makeatletter\def\@captype{table}
% \vspace{-0.05mm}
% \caption{\textbf{Performance of VLMs in NExT-QA and FunQA.} N.QA, F.MC, and F.MC-R denote NExT-QA, FunQA-MC and FunQA-MC-R, which are multi-choice datasets and we use \textbf{accuracy} as the metric. Regarding N.QE (NExT-OE) and FunQA, which are open-ended answer datasets, we employ metrics based on \textbf{GPT-4}. $^{\ast}$ indicates that this model utilizes a Chat-box mechanism, and we employ GPT-3.5 to automatically convert its output into multiple-choice answers; $^{\dagger}$ signifies the GPT-4 model without any visual information input; $^{\ddagger}$ denotes the GPT-4V(ision) model, where an average sample of four frames of images is used as input. $^{\mathsection}$ denotes an sample version ($10\%$) of all free-text QA pairs.}
% \label{T:vlm_next_funqa}
% \vspace{-5pt}
% \centering
% \setlength\tabcolsep{10pt}
% \resizebox{1\linewidth}{!}{
% \begin{tabular}{lcccp{1cm}<{\centering}cc}
% \toprule
% Metrics & \multicolumn{3}{c}{Acc} & \multicolumn{2}{c}{GPT-4} \\
% \cmidrule{2-4}\cmidrule{5-6}
% Dataset & N.QA & F.MC & F.MC-R & N.QE & FunQA $^{\mathsection}$ \\
% \midrule
% Random & 20 & 20 & 20 & - & - \\
% Otter (D.C.) $^{\ast}$ & 35 & 27 & 17 & 54 & 22 \\
% Otter (FunQA) $^{\ast}$ & 42 & 31 & 26 & 58 & \textbf{28} \\
% GPT-4 $^{\dagger}$ & 44 & 34 & 23 & 30 & 2 \\
% GPT-4V $^{\ddagger}$ & \textbf{80} & \textbf{61} & \textbf{39} & \textbf{79} & 5\\
% \bottomrule
% \end{tabular}
% }
% \end{minipage}
% \end{minipage}
% \end{center}

% For instance, when the Ground Truth answer is \textcolor{blue}{'dog'} and HGA responds with \textcolor{red}{'cat'} indicating an incorrect answer, WUPS erroneously assigns a high score of \textcolor{red}{85.6} due to both being animals. Conversely, Otter provides the answer \textcolor{green}{'A \textbf{dog} is shown laying on the floor in front of the toddlers'}, which is entirely accurate. However, it receives an excessively low score of \textcolor{green}{0.00015} due to long-word count. In contrast, the GPT-4-based evaluation metrics assign scores of \textcolor{red}{0} and \textcolor{green}{80} to the respective responses. 
% Specific examples can be found in Appendix \ref{appendix_d: Exp_result}. 
% \begin{table}[h]
\vspace{-20pt}
\caption{\textbf{Model's Performance in NExT-OE.} This table illustrates that while the previous SOTA model, HGA, excels in the classic WUPS metric, it underperforms with GPT-4. See Appendix~\ref{appendix_d: Evaluation Metrics} for a discussion on potential issues with WUPS.}
\vspace{-10pt}
\label{T:vlm_nextoe}
\centering
\setlength\tabcolsep{5pt}
\resizebox{1.0\linewidth}{!}{
\begin{tabular}{lp{5cm}<{\centering}p{5cm}<{\centering}}
\toprule
Metrics & WUPS & GPT-4  \\
% \toprule
\hline
HGA (SOTA on NExT-OE)& 25.18	& 25.06 \\
Otter (D.C.) & 1.26 & 64.89 \\
Otter (FunQA) & 0.79 & \textbf{73.06} \\
\bottomrule
\end{tabular}
}
\vspace{-20pt}
\end{table}



% These findings demonstrate: 1) Traditional evaluation metrics for VLM-generated free-text responses are no longer meaningful, while GPT-4-based metrics offer reasonable assessments. 2) Otter excel at addressing NExT-OE. 3) The superior performance of Otter (FunQA) in comparison to Otter (DC) serves as evidence of the model's enhanced capabilities in visual understanding and reasoning following fine-tuning on the FunQA dataset.

\noindent\textbf{NExT-QA is not a challenge for GPT-4.}\quad
Our investigation into VLMs' performance on NExT-QA and FunQA datasets, shown in Table \ref{T:vlm_next_funqa}, reveals key insights. In the multi-choice format, GPT-4, even without video frames, scores comparably to VLMs on NExT-QA but resembles random guessing on FunQA-MC. Furthermore, while GPT-4V scores 80 and 61 on NExT-QA and FunQA-MC respectively, it drops to 39 on FunQA-MC-R, which focuses on counter-intuitive reasoning. This suggests that NExT-QA's inferential questions are no longer challenging for GPT-4, and the significant score difference in FunQA variants underlines GPT-4V's struggle with complex, non-intuitive questions.
% \begin{figure}[b]
%     \centering
%     \includegraphics[width=1.0\linewidth]{figures/vlm_next_fun.png}
%     \vspace{-0.5cm}
%     \caption{\textbf{Error rate of VLMs in NExT-QA and FunQA}}
%     \label{fig:vlm_next_funqa}
%     \vspace{-0.4cm}
% \end{figure}

\noindent\textbf{FunQA is challenging for VLMs and GPT-4V.}\quad
As Table~\ref{T:vlm_next_funqa} indicates, GPT-4 scores well on NExT-OE even without video access. The contrast in its performance on NExT-OE versus FunQA highlights FunQA's complexity in video reasoning. A significant performance disparity between GPT-4 and its video-enhanced version, GPT-4V, on FunQA again underscores the value of video content, corroborating findings from Table \ref{T:onlyframes}. Overall, FunQA stands out from prior benchmarks with its focus on reasoning skills and high-quality QA pairs deeply linked to video content, establishing itself as a robust benchmark in the LLM era for assessing VLMs' capabilities in counter-intuitive reasoning.

\begin{table}[t]
\vspace{-0.05mm}
\caption{\textbf{Performance of VLMs in NExT-QA and FunQA.} N.QA, F.MC, and F.MC-R denote NExT-QA, FunQA-MC and FunQA-MC-R, which are multi-choice datasets and we use \textbf{accuracy} as the metric. Regarding N.QE (NExT-OE) and FunQA, which are open-ended answer datasets, we employ metrics based on \textbf{GPT-4}. $^{\ast}$ indicates that this model utilizes a Chat-box mechanism, and we employ GPT-3.5 to automatically convert its output into multiple-choice answers; $^{\dagger}$ signifies the GPT-4 model without any visual information input; $^{\ddagger}$ denotes the GPT-4V(ision) model, where an average sample of four frames of images is used as input. $^{\mathsection}$ denotes an sample version ($10\%$) of all free-text QA pairs.}
\label{T:vlm_next_funqa}
\vspace{-10pt}
\centering
\setlength\tabcolsep{10pt}
\resizebox{1.0\linewidth}{!}{
\begin{tabular}{lcccp{1cm}<{\centering}cc}
\toprule
Metrics & \multicolumn{3}{c}{Acc} & \multicolumn{2}{c}{GPT-4} \\
\cmidrule{2-4}\cmidrule{5-6}
Dataset & N.QA & F.MC & F.MC-R & N.QE & FunQA $^{\mathsection}$ \\
\midrule
Random & 20 & 20 & 20 & - & - \\
Otter (D.C.) $^{\ast}$ & 35 & 27 & 17 & 54 & 22 \\
Otter (FunQA) $^{\ast}$ & 42 & 31 & 26 & 58 & \textbf{28} \\
GPT-4 $^{\dagger}$ & 44 & 34 & 23 & 30 & 2 \\
GPT-4V $^{\ddagger}$ & \textbf{80} & \textbf{61} & \textbf{39} & \textbf{79} & 5\\
\bottomrule
\end{tabular}
}
\vspace{-20pt}
\end{table}
 

%several reasons: 1) It comprises high-quality and focuses on tasks that emphasize a model's reasoning abilities. 2) The questions are closely intertwined with the video content. 3) It provides high-quality free-text annotations. Collectively, these attributes position FunQA as a consistent benchmark in the era of Large Language Models (LLMs) for evaluating VLMs, particularly in the domain of counter-intuitive reasoning abilities.



%it can be concluded that traditional metrics appear less suitable for more in-depth evaluation.

%The inadequate and unbalanced information and the lagging loos function may lead to ineffective training.

